%% Part 1 — Why completeness proofs need a machine (paper §1).

\PaperPart
\section{Why completeness proofs need a machine}

\begin{frame}{Circuits are programs; rules are their equational theory \InPaper}
\begin{columns}[c]
\begin{column}{0.55\textwidth}
A circuit optimizer, a compiler pass, a diagram normalizer: the
justification is always the same --- \alert{local rewrite rules} on a
syntactic object, preserving a \alert{semantics}.
\vspace{6pt}
\begin{center}
\begin{quantikz}[row sep=0.2cm,column sep=0.18cm,font=\scriptsize]
  & \gate{H} & \gate{H} & \gate{S} & \ctrl{1} & \\
  & & & & \control{} &
\end{quantikz}
$\;\overset{\scriptscriptstyle H^2=\varepsilon}{\approx}\;$
\begin{quantikz}[row sep=0.2cm,column sep=0.18cm,font=\scriptsize]
  & \gate{S} & \ctrl{1} & \\
  & & \control{} &
\end{quantikz}
$\;\overset{\scriptscriptstyle\text{diagonal}}{\approx}\;$
\begin{quantikz}[row sep=0.2cm,column sep=0.18cm,font=\scriptsize]
  & \ctrl{1} & \gate{S} & \\
  & \control{} & &
\end{quantikz}
\end{center}
\vspace{4pt}
Two questions attach to every rule set:
\begin{itemize}
  \item \textbf{Soundness} --- does every rule preserve the semantics?\\
        {\small\color{gray}finitely many small algebraic identities: easy.}
  \item \textbf{Completeness} --- is every true equation derivable?\\
        {\small\color{gray}a different animal altogether.}
\end{itemize}
\end{column}
\begin{column}{0.42\textwidth}
\centering
\begin{tikzpicture}[
  >={Stealth[length=2mm]}]
  \node[blob=accent] (S) at (0,0) {};
  \node[blob=paperC] (M) at (3.3,0) {};
  \node[font=\small\bfseries,accent] at (0,2.0) {syntax};
  \node[font=\small\bfseries,paperC] at (3.3,2.0) {semantics};
  \node[font=\scriptsize,accent] at (0,1.15) {circuits / $\approx$};
  \node[font=\scriptsize,paperC] at (3.3,1.15) {unitaries, \dots};
  \node[circle,fill=accent,inner sep=1.3pt,label=left:{\small $w$}] (w) at (-0.2,0.35) {};
  \node[circle,fill=accent,inner sep=1.3pt,label=left:{\small $v$}] (v) at (-0.2,-0.55) {};
  \node[circle,fill=paperC,inner sep=1.3pt,label=right:{\small $U$}] (u) at (3.2,-0.1) {};
  \draw[->,thick,gray] (w) to[bend left=12] node[above,font=\scriptsize]{$\llbracket\_\rrbracket$} (u);
  \draw[->,thick,gray] (v) to[bend right=12] (u);
  \draw[<->,thick,newC,dashed] (w) to[bend right=40] node[left,font=\scriptsize,align=right]{derivable\\from rules?} (v);
\end{tikzpicture}

\small soundness: $w\approx v\;\Rightarrow\;\llbracket w\rrbracket=\llbracket v\rrbracket$\\
completeness: $\llbracket w\rrbracket=\llbracket v\rrbracket\;\Rightarrow\;w\approx v$
\end{column}
\end{columns}
\end{frame}

\begin{frame}{How completeness is proved on paper \InPaper}
\begin{columns}[c]
\begin{column}{0.47\textwidth}
The standard strategy has three steps:
\begin{enumerate}
  \setlength{\itemsep}{4pt}
  \item define a \alert{normal form} for circuits;
  \item show every circuit \alert{rewrites to normal form} using the rules;
  \item show \alert{distinct normal forms denote distinct} semantic objects.
\end{enumerate}
\vspace{6pt}
Step 2 is where it explodes: for \emph{every} rule and \emph{every} way a
generator can meet a normal form, the combination must reduce again ---
mechanical, unforgiving, enormous.
\end{column}
\begin{column}{0.50\textwidth}
\centering
\begin{tikzpicture}[x=3.3mm,y=3.3mm]
  % a grid of "cases": generators x normal-form shapes
  \foreach \i in {0,...,13} {
    \foreach \j in {0,...,8} {
      \pgfmathsetmacro{\shade}{mod(\i*7+\j*13,11)}
      \ifdim\shade pt<1pt
        \fill[newC!55] (\i,\j) rectangle ++(0.85,0.85);
      \else
        \fill[paperC!35] (\i,\j) rectangle ++(0.85,0.85);
      \fi
    }
  }
  \node[rotate=90,font=\scriptsize,anchor=south] at (-0.3,4.4) {generator $g$ (incl.\ every wire)};
  \node[font=\scriptsize,anchor=north] at (7,-0.2) {shape of the normal form it meets};
  \node[font=\scriptsize,anchor=south,align=center] at (7,9.1)
     {one case per (rule, coset) or (generator, coset)};
  \fill[paperC!35] (15,6) rectangle ++(0.85,0.85); \node[font=\scriptsize,anchor=west] at (16,6.4) {routine};
  \fill[newC!55] (15,4) rectangle ++(0.85,0.85); \node[font=\scriptsize,anchor=west,align=left] at (16,4.4) {a real\\derivation};
\end{tikzpicture}
\end{column}
\end{columns}
\end{frame}

\begin{frame}{The literature wears the scars \InPaper}
\begin{columns}[c]
\begin{column}{0.50\textwidth}
\small
\begin{itemize}
  \setlength{\itemsep}{3pt}
  \item Selinger (2015): $n$-qubit Clifford --- a substantial rewrite
        system and normal-form analysis.
  \item Li, Mosca, Ross, van de Wetering, Zhao (2025): qutrit Clifford, first
        complete fragment in odd prime dimension --- \alert{$>50$ pages} of proofs.
  \item Clément, Heurtel, Mansfield, Perdrix, Valiron (2023, 2024):
        complete theory for full quantum circuits, and its minimal refinement.
  \item Clément (LICS 2026): Real-Clifford+$CH$ --- a 13-page paper with an
        \alert{appendix of 264 pages} of equational derivations.
  \item Rig-categorical completeness (POPL '22, '24, '26); Blake (2026)
        simplifies the rule sets of six near-Clifford fragments.
\end{itemize}
\end{column}
\begin{column}{0.47\textwidth}
\centering
\begin{tikzpicture}
\begin{axis}[xbar, width=5.0cm, height=4.6cm, xmin=0, xmax=300,
  symbolic y coords={Real-Clifford+CH appendix,Real-Clifford+CH paper,Qutrit Clifford proofs},
  ytick=data, y tick label style={font=\scriptsize,align=right,text width=2.1cm},
  x tick label style={font=\scriptsize}, xlabel={\scriptsize pages},
  nodes near coords, nodes near coords style={font=\scriptsize},
  bar width=9pt, enlarge y limits=0.3, axis lines*=left]
\addplot[fill=newC!70,draw=newC] coordinates
  {(264,Real-Clifford+CH appendix) (13,Real-Clifford+CH paper) (50,Qutrit Clifford proofs)};
\end{axis}
\end{tikzpicture}

{\small Long, regular, mechanical: exactly the proofs a proof assistant
should check --- it does not tire of case splits, and where a case is a
computation, its \alert{evaluator} can carry it out.}
\end{column}
\end{columns}
\end{frame}

\begin{frame}{Why this is a programming-languages problem \InPaper}
\begin{columns}[T]
\begin{column}{0.50\textwidth}
\begin{itemize}
  \setlength{\itemsep}{5pt}
  \item Circuit calculi \emph{are} programming languages; completeness of
        their equational theory is a PL question.
  \item It licenses \alert{normal-form decision procedures}, and tells an
        optimizer author that \alert{no valid rewrite is missing}.
  \item It underlies the rig-categorical completeness theorems for
        reversible and quantum languages (Choudhury--Karwowski--Sabry POPL'22,
        Carette--Heunen--Kaarsgaard--Sabry POPL'24, Fang--Heunen--Kaarsgaard
        POPL'26).
  \item The contribution is a \alert{design}: objects that both
        \emph{compute} and \emph{prove}.
\end{itemize}
\end{column}
\begin{column}{0.46\textwidth}
\begin{keybox}[title={Announced up front: setoids, not quotients}]
\small
Plain, setoid-based Agda on the standard library --- not cubical, not HoTT.
A presented monoid is a \emph{setoid on words}, not a quotient type.
\begin{itemize}
  \item the library \alert{computes with syntax}: coset tables and normal
        forms pattern-match on concrete representatives, obligations close
        by evaluation;
  \item central devices are maps \alert{into} the syntax (sections, coset
        representatives, amalgamation tails);
  \item derivations are \alert{data} one can induct on.
\end{itemize}
\end{keybox}
\end{column}
\end{columns}
\end{frame}

\begin{frame}{Contributions of the paper \InPaper}
\small
\begin{columns}[T]
\begin{column}{0.50\textwidth}
\begin{enumerate}
  \setcounter{enumi}{-1}
  \setlength{\itemsep}{4pt}
  \item \textbf{A library} for circuit normal forms and complete relation
        sets: presentations as inductive congruences; normal forms as stdlib
        function bundles; generic \aid{by-normalization}; a verified
        Reidemeister--Schreier engine; solvers by evaluation.
  \item \textbf{Inductively defined circuits}: wire-indexed generators,
        structural rules supplied once for every gate set.
  \item \textbf{A catalogue} of verified presentations: $\mathbb{Z}_n$,
        $S_n$, $\mathbb{Z}_N\wr S_n$, Pauli $(\mathbb{Z}_p{\times}\mathbb{Z}_p)^n$,
        qubit/qutrit Clifford+$T$, $U_3(\mathbb{Z}[\frac12,i])$.
\end{enumerate}
\end{column}
\begin{column}{0.48\textwidth}
\begin{enumerate}
  \setcounter{enumi}{2}
  \setlength{\itemsep}{4pt}
  \item \textbf{Monoid-level} amalgamation and Reidemeister--Schreier
        (after Bian--Selinger).
  \item \textbf{Presentation theorems} for direct, semidirect, $n$-fold and
        amalgamated products, and group extensions.
  \item \textbf{New stdlib theorems}: semidirect products, amalgamated free
        products of monoids, extensions, epimorphisms, $\mathsf{Perm}(n)$ as a group.
\end{enumerate}
\vspace{4pt}
\begin{paperbox}
Every headline theorem is restated with its full type in
\texttt{MainTheorems.agda}; \texttt{--safe}, \texttt{--cubical-compatible},
no postulates, no holes, no classical axioms.
\end{paperbox}
\end{column}
\end{columns}
\end{frame}
